\section{Numerical studies}
\label{sec:numerical-comparison}

In this section we carry out some numerical studies to give some insight into how well randomized Trotter extrapolation (or Trotter extrapolation in general) may perform in practice.

We start by explicitly plotting the error bounds we derive for an electronic structure problem. We plot both bounds using the LKW extrapolation schedule, as well as extrapolation schedules found by simple brute-force search. We see that under the assumptions made, randomized Trotter extrapolation starts to beat regular Trotter (in upper bounds) at a mild target precision. Moreover, extrapolation schedules generated by simple brute-force search improve on the asymptomatically well-conditioned LKW extrapolation schedule. We detail the

\subsection{Constant factor comparison}\label{sec:trotter-constants}

We analyze resource costs for the Hamiltonian electronic structure of plane-waves, for which it is known $\acomm^{(j)} = O(n^j)$, where $n$ is the number of basis functions. For simplicity, we consider the task of compiling a time signal. In order to generate resource estimates, we pull out constant factors for randomized Trotter extrapolation and regular Trotterized time evolution.

\medskip

\textbf{Constant factors for regular product formulae.} The number of Trotter steps for a $p$th-order product formula to attain additive error $\varepsilon$ in the time signal is shown in \cite[Eq.~(189)]{childs2021TheoryTrotter} to be bounded as
\begin{equation}\label{eq:trotter-steps}
r_{\mathrm{Trotter}} \leq \left(\frac{2}{1+p}\right)^{\frac{1}{p}}
\bigl(\alpha_{\mathrm{comm}}^{(p+1)}\bigr)^{\frac{1}{p}}
(\Upsilon T)^{1+\frac{1}{p}}
\varepsilon^{-\frac{1}{p}}
.
\end{equation}
For our resource comparison, we take the Suzuki-Trotter formulae \cite{suzuki1991general}, for which the number of stages for the $2k$th order formula (for $k\geq 1$) is
\begin{equation}
    \Upsilon_{2k} = 2 \cdot 5^{k-1}
\end{equation}
Consequently, the total gate depth for a Suzuki-Trotter formula is 
\begin{align}
    \Gamma
\bigl(\alpha_{\mathrm{comm}}^{(p+1)}\bigr)^{\frac{1}{p}}
T^{1+\frac{1}{p}}
\varepsilon^{-\frac{1}{p}} \quad &\text{for $p=1$}\\
\Gamma \left(\frac{2}{1+p}\right)^{\frac{1}{p}} \cdot \left( 2 \cdot 5^{k-1} \right)^{2+\frac{1}{p}}
\bigl(\alpha_{\mathrm{comm}}^{(p+1)}\bigr)^{\frac{1}{p}}
T^{1+\frac{1}{p}}
\varepsilon^{-\frac{1}{p}} \quad &\text{for $p=2k$; $k \geq 1$}
\end{align}

\medskips

\textbf{Refined condition number.} We make a further refinement to the bounds in Lemmas \ref{lem:generic-Richardson-error-2} and \ref{lem:specialized-Richardson-error}. We recall in both derivations we  amplify the error series remainder by the full extrapolation coefficient norm $\|\vec{b}\|_1$. However this is an overestimate. Recall from Lemma~\ref{lem:time-evol-error-series} that each Trotterized time evolution operator admits the operator-valued expansion
\begin{equation}\label{eq:trotter-series-recall}
\cP^{1/s_k}(s_k T) - e^{-iHT} = \sum_{\substack{j \in \sigma \mathbb{Z}_+\\
j \geq p}} s_{k}^{\,j}\, \tilde{E}_{j+1}(T)\,,
\end{equation} with $s_k =  s/ q_k$. Inserting this into the extrapolation $\cP^{(R)}_{p, m}(T) =  \sum_{k=1}^m b_k \cP^{1/s_k} (s_k T)$ and using the fact that the schedule $\{b_k, q_k\}_{k=1}^m$ cancels the powers $s^{p}, s^{p+\sigma}, \ldots, s^{p + \sigma (m-1)}$ (as in Lemma~\ref{lem:specialized-Richardson-error}), the residual error operator takes the form
\begin{align}\label{eq:tilde-b-bound}
\cP^{(R)}_{p, m}(T)- e^{-iHT}\;&=\;\sum_{\substack{j \in \sigma\mathbb{Z}_+ \\ j \geq \sigma(m-1)+p}} s^{\,j}\ \Big(\sum_{k=1}^{m} \frac{b_k}{q_k^{\,j}}\Big) \tilde{E}_{j+1}(T) \\
&= \;\sum_{\substack{j \in \sigma\mathbb{Z}_+ \\ j \geq \sigma(m-1)+p}} s_1^{\,j}\ \sum_{k=1}^{m} {b_k} \Big(\frac{q_1}{q_k}\Big)^{j} \tilde{E}_{j+1}(T) \\
&\leq \; \sum_{k=1}^{m} {b_k} \Big(\frac{q_1}{q_k}\Big)^{p} \sum_{\substack{j \in \sigma\mathbb{Z}_+ \\ j \geq \sigma(m-1)+p}} s_1^{\,j} \tilde{E}_{j+1}(T) \,.
\end{align}
As a result, for Lemma \ref{lem:specialized-Richardson-error}, rather than $\|\vec{b}\|_1$, the remainder can be thought of as being amplified by $\|\vec{b}^{(p)}\|_1$, the norm of the vector whose entries are $b_kq_1^p/q_k^p$. Similarly, for Lemma \ref{lem:generic-Richardson-error-2} the remainder is amplified by $\|\vec{b}^{(1)}\|_1$. In practice this can lead to much smaller bounds when few extrapolation points are used. We note however that the sample complexity is still controlled by $\|\vec{b}\|_1$.

\medskip

\textbf{Constant factors for plane-wave dual basis.} For electronic-structure Hamiltonians in the plane-wave dual basis with $n$ basis functions it is known that 
\begin{equation}\label{eq:comm-scaling-basic}
\alpha^{(j)}_{\text{comm}} = O(n^{j}),
\qquad
\lambda_{\text{comm}} = O(n),
\end{equation}
from \cite{childs2021TheoryTrotter} and \cite{chakraborty2025quantum} respectively. Thus, asymptotically, the dependence of extrapolated product formulae on $n$ is the same as that for standalone product formulae. In the following Lemma we prove a result with refined constant factors for any exponentially growing $\acomm^{(j)}$.

\begin{lemma}[Bound on $\lambda_{\mathrm{comm}}$ for exponentially growing $\acomm^{(j)}$]\label{lem:lcomm-tight-bound}
Consider a $p$th order product formula and let $\alpha_{\mathrm{comm}}^{(j)}$ and $\lcomm$ be defined as in Definitions \ref{def:acomm} and \ref{def:lcomm} respectively.
Assume there exists $A > 0$ such that
\begin{equation}\label{eq:alpha-comm-bound}
\alpha_{\mathrm{comm}}^{(t)} \le A\, n^{j}, \qquad \forall\, j \ge 1.
\end{equation}
Then, we have
\begin{equation}\label{eq:lambda-comm-bound-explicit}
\lambda_{\mathrm{comm}}
\;\le\;
n \cdot \max\{A^{\frac{1}{p+1}}, 1\} \cdot \sup_{0\leq k \leq 1/p} \left(\frac{e }{(p + 1)^{2}}\left( \frac{1}{k} - p \right) \right)^{k}\,.
\end{equation}
\end{lemma}


\begin{proof}
    Starting from Definition \ref{def:lcomm} and imposing the bound in Eq.~\eqref{eq:alpha-comm-bound}, we have
    \begin{equation}
    \prod_{\kappa=1}^{\ell}
    \frac{\alpha_{\mathrm{comm}}^{(j_\kappa + 1)}}{(j_\kappa + 1)^2}
    \le
    A^{\ell} n^{j+\ell}
    \prod_{\kappa=1}^{\ell} (j_\kappa + 1)^{-2}.
    \end{equation}
    Since $j_\kappa + 1 \ge p + 1$, we have $\prod_{\kappa} (j_\kappa + 1)^{-2} \le (p + 1)^{-2\ell}$.
    The number of integer tuples $\{j_1, \ldots, j_\ell\}$ satisfying $\sum_{\kappa} j_\kappa = j$ with $j_\kappa \ge p$ is bounded by
    \begin{equation}
    \sum_{\substack{j_1, \dots, j_\ell \ge p \\ j_1 + \cdots + j_\ell = j}} 1
    =
    \binom{j - \ell p + \ell - 1}{\ell - 1}
    \le
    \left(\frac{e (j-\ell p+\ell-1)}{\ell-1}\right)^{\ell-1}.
    \end{equation}
    by the standard bound on binomial coefficients.
    Combining these bounds yields
    \begin{equation}
    \sum_{\substack{j_1, \dots, j_\ell \ge p \\ j_1 + \cdots + j_\ell = j}}
    \prod_{\kappa=1}^{\ell}
    \frac{\alpha_{\mathrm{comm}}^{(j_\kappa + 1)}}{(j_\kappa + 1)^2}
    \le n^{j+\ell}
    \frac{A^{\ell} }{(p + 1)^{2\ell}}   \left(\frac{e (j-\ell p+\ell-1)}{\ell-1}\right)^{\ell-1} \leq n^{j+\ell}
    A^{\ell}    \left(\frac{e }{(p + 1)^{2}} \frac{  j-\ell p+\ell-1}{\ell-1}\right)^{\ell-1}\,.
    \end{equation}
    Taking the $(j+\ell)$-th root and the supremum, we can write
    \begin{equation}
        \lcomm \leq n \sup_{\substack{j \in \mathbb{Z}^+_{\ge \sigma m} \\ 1 \le \ell \le \lfloor j/p \rfloor}} \left( A^{\frac{\ell}{j + \ell}} \right) \sup_{\substack{j \in \mathbb{Z}^+_{\ge \sigma m} \\ 1 \le \ell \le \lfloor j/p \rfloor}} \left(\left(\frac{e }{(p + 1)^{2}} \frac{j-\ell p+\ell-1}{\ell-1}\right)^{\frac{\ell-1}{j+ \ell}} \right)\,.
    \end{equation}
    The first supremum takes value $A^{\frac{1}{p+1}}$ for $A \geq 1$.
    The second supremum can be bounded as
    \begin{align}
        \sup_{\substack{j \in \mathbb{Z}^+_{\ge \sigma m} \\ 1 \le \ell \le \lfloor j/p \rfloor}} \left( \left( \frac{e }{(p + 1)^{2}} \frac{j-\ell p+\ell-1}{\ell-1}\right)^{\frac{\ell-1}{j+ \ell}} \right) &\leq  \sup_{0\leq k \leq 1/p} \inf_{\ell; \geq 0} \left(\frac{e }{(p + 1)^{2}} \left( \frac{1}{k} - \frac{\ell p +1}{\ell-1} \right) \right)^{k}  \\
        &\leq  \sup_{0\leq k \leq 1/p} \left(\frac{e }{(p + 1)^{2}}
         \left( \frac{1}{k} - p \right) \right)^{k}\,,
    \end{align}
    where $\inf_{\ell; ()\geq 0}$ indicates an infimum such that the quantity inside the brackets is non-negative. 
\end{proof}

When $A \geq 1$, the ratio of $\lcomm^{1+1/p}$ to $(\acomm^{(p+1)})^{1/p}$ (which represents respective contributions to gate complexities) is
\begin{equation}
    \frac{\lcomm^{1+1/p}}{(\acomm^{(p+1)})^{1/p}} = \sup_{0\leq k \leq 1/p} \left(\frac{e }{(p + 1)^{2}}\left( \frac{1}{k} - p \right) \right)^{k}\,.
\end{equation}
We can numerically determine this value and find it to be $\leq 1.5035$ for $p=1$, $\leq 1.1487$ for $p=2$ and $\leq 1.0445$ for $p=4$. Assuming $A\geq 1$, the choice $A=1$ maximizes the bound for $\lcomm$ (it is the worst-case choice). From hereon we take this choice.

\medskip

\textbf{Choice of extrapolation schedule.} In Lemmas \ref{lem:generic-Richardson-error-2} and \ref{lem:specialized-Richardson-error} we demonstrate bounds on the maximum number of Trotter steps for two types of extrapolation schedules. The first extrapolation schedule removes error series powers starting from $s^1$ (or $s^2$ for symmetric error series), and is compatible with the well-conditioned extrapolation schedule of \cite{low2019Multiproduct} which we used to derive our main theoretical results (Theorems \ref{thm:overlap}, \ref{thm:overlap-with-observable}, \ref{thm:distribution-recovery}). The second uses the information that a $p$th order Trotterized time evolution error series has leading order power $s^p$. In this setting we will find find extrapolation schedules by simple brute-force optimization. In both settings, we take Trotter-Suzuki product formulae where we can use the bound $a_{\max} \leq 1$.

We note that in both bounds there is a ceiling function, and we omit this for simplicity -- this would be valid if $\Upsilon \lcomm T^{1+1/p} \log(1/\varepsilon)$ were large compared to the smallest error amplification factor $q_1$. This simplification allows us to compare bounds with regular Trotter formulae (Eq.~\eqref{eq:trotter-steps}) such that we can disregard the dependency on $\Upsilon$ and $T$, as the dependence of all bounds with respect to these parameters is the same.

For the LKW well-conditioned grid we plot the following quantity:
\begin{equation}
    \frac{q_{\max}}{q_{\min}} \lambda_{\text{scale}}^{1+1/p} \left(\frac{a(\varepsilon)\|\vec{b}^{(1)}\|_1}{\varepsilon}\right)^{\frac{1}{\sigma m}}\,,
\end{equation}
where (as will all schemes) we omit dependencies on $\Upsilon$, $T$ and $n$ due to identical contribution -- we denote $\lambda_{\text{scale}}$ as the upper bound in Lemma \ref{lem:lcomm-tight-bound} divided by $n$, i.e.~the constant factor contribution to the extrapolated commutator factor. The quantity $a(\varepsilon)$ was chosen to be $4$ for simplicity in our previous results following the analysis in Lemma \ref{lem:generic-Richardson-error} -- here, we numerically pick a refined constant for each target precision chosen to be the smallest number such that an upper bound on step number condition $s_1 (a_{\max } \Upsilon \lcomm T)^{1+1/p}$ is consistent. 

For our brute-force search we optimize plot the following objective over a simple search space of all possible integer sequences $\{q_k\}_{k=1}^m$
\begin{equation}\label{eq:optimization-objective}
\mathcal{C}(q_1, \ldots, q_m) = \frac{q_{\max}}{q_{\min}} \lambda_{\text{scale}}^{1+1/p} \left(\frac{a(\varepsilon)\|\vec{b}^{(p)}\|_1}{\varepsilon}\right)^{\frac{1}{\sigma (m-1)+p}},
\end{equation}
in the domain $q_k \in [1, 10]$, and where $\vec{b}$ is computed from as the solution of the Vandermonde linear system as in Eq.~\eqref{eq:vandermonde-system}. 

We present our plot of upper bounds of the LKW well-conditioned grid, the brute-force optimized grid, alongside the regular Trotter bound across orders $p \in {1,2,4}$ in Figure~\ref{fig:overhead-plots}. We recall that the LKW grid only works for symmetric product formulae, and so is not avaialbe for $p=1$. We record the sample overhead amplification factor $\|\vec{b}\|_{1}^{2}$ as the color of each data marker. During brute-force search, candidate grids with $\|\vec{b}\|_{1}^{2} > \texttt{thresh}$ are rejected for $\texttt{thresh} \in  \{10,100,1000\}$, leading to three distinct lines. 

\begin{figure}[ht]
    \centering
    \includegraphics[width=\linewidth]{images/overhead-multi.pdf}
    \caption{\textbf{Bounds on maximum number of Trotter steps.} We plot our best known bounds on the maximum required Trotter step number to attain a given additive error for the time signal task. The number we plot at each target precision is proportional to the true Trotter step number, where the true value contains a factor dependent on the target simulation time and problem size.}
    \label{fig:overhead-plots}
\end{figure}

We observe first that despite adopting a very simplistic search strategy for extrapolation schedules, our brute-force optimization improves upon the prescribed extrapolation schedule of~\cite{low2019Multiproduct}, even whilst asking a modest increase in sample budget of $\leq 10$. This demonstrates that heuristic optimization could lead to useful extrapolation schedules that go beyond the performance of our (asymptotically strong) analytical results.  Second, our bounds on extrapolated Trotter formulae outperform Trotter formulae for small enough precision, but is not found for large precision. We remark that for each order the extrapolation schemes at the largest target precision are the trivial one (no extrapolation) -- thus, the deviation between extrapolated data and the regular Trotter curve is purely due to difference in constant factors -- were constant factors to be equal, randomized Trotter extrapolation would outperform regular Trotter unconditionally across our studied orders. %This is seen most explicitly at error $\varepsilon=1$, where the trivial solution is found to be optimal for $p=1$ and $p=2$, and the deviation in results arises purely due to different constant factors. 
We thus envisage improvements can be found with refined analysis of constant factors and a refined optimization procedure in future work.

\medskip

\textbf{Gate depths.} In the previous figure we have only investigated Trotter step number, which is only one part of the contribution to total gate complexity per circuit. Inspection of Figure \ref{fig:overhead-plots} shows that the crossover point where randomized Trotter extrapolation is more efficient than regular Trotter appears to get worse as the order increases. In reality, the number of stages prohibits use of high-order product formulae (e.g.~the sixth-order Trotter--Suzuki formula requires $50$ operators per Hamiltonian term per Trotter step, compared to $2$ for second-order). Figure~\ref{fig:gate-depth} presents (a number proportional to) the gate depth for the same data displayed in Figure \ref{fig:overhead-plots}. For simplicity, for lines corresponding to extrapolated data we have selected only extrapolation schedules found via heuristic search with sample amplification factor at most $10$, and recall that stronger performance can be obtained if a larger sample budget is allowed. We choose a system size of $n=100$, deemed a large enough regime to observe potential quantum advantage for the uniform electron gas (Jellium) \cite{babbush2018LowDepthQSimMaterial}. In the setting chosen for comparison we see that randomized Trotter extrapolation beats regular product formulae (in bounds) across order $p=1$ to $p=6$ below any error rate $\leq 2\times 10^{-2}$.

 \begin{figure}[ht]
    \centering
\includegraphics[width=\linewidth]{images/gate_depth.pdf}
    \caption{\textbf{Bounds on gate depth.} We plot best available bounds on (a quantity proportional to) the gate depth for a plane-wave dual basis Hamiltonian with $100$ basis functions. \textit{Left:} Superposed data taken from Figure \ref{fig:overhead-plots} where we have selected brute-force optimized extrapolation schedules with sample overhead factor 10. \textit{Right:} We highlight the best curve across all orders. The extrapolated bound beats the Trotter bound at $\sim 3 \times 10^{-3}$}
    \label{fig:gate-depth}
\end{figure}


\medskip

\textbf{Exact error.} So far we have compared analytical bounds on the gate depth to reach some target precision. This is useful if one wishes to perform resource estimates, or to perform an experiment with precision guarantees (assuming perfect execution). However, bounds only tell one side of the story, and the true error attained by randomized Trotter extrapolation may be much better than the bounds we derive in practice. Indeed, the true error of regular Trotter circuits may also be substantially better. 

In order to gain some insight into the exact error attainable by randomized Trotter extrapolation, we take the extrapolation schedule data from our previous numerical studies and import them in a "blind" fashion into a new problem: simulating a Heisenberg Hamiltonian. We consider an anisotropic Heisenberg Hamiltonian with an external magnetic field
\begin{equation}
    H = J \sum_{i=1}^{L-1} ( X_i X_{i+1} + Y_i Y_{i+1} + 0.9 Z_i Z_{i+1}) - h\sum_{i=1}^L Z_i\,,  
\end{equation}
with a choice of $J=h=1$.
On a system of 8 qubits we evaluate the exact operator norm error of a simulation of the time evolution operator for both a Suzuki-Trotter circuit and randomized Trotter extrapolation implementations for $p=2$. We import all of our extrapolation schemes found for $p=2$ in our previous numerical study. We display our results in the left subfigure of Figure \ref{fig:exact}. We see that anything below a very large operator norm error ($\sim 1$), all the previously found extrapolation schemes allow substantially reduced error rates compared to regular Trotter circuits.

In the right subfigure of Figure \ref{fig:exact} we investigate the robustness of our extrapolation schemes in the presence of noise. We consider a simple model of measurement noise described by random Pauli-$X$ rotations. This incurs a noise floor below which no strategy can pass. We see that the extrapolated data is able to hit the noise floor using fewer Trotter steps than the regular Trotter circuits. We see evidence that heuristically found extrapolation schemes are in practice robust to noise, despite having non-trivial condition numbers.

\begin{figure}[ht]
    \centering
    \includegraphics[width=\linewidth]{images/exact-error.pdf}
    \caption{\textbf{Exact error for Heisenberg model.} We plot the exact error on the time evolution operator for a Heisenberg chain Hamiltonian on $8$ qubits using the second-order Suzuki-Trotter formula. Brute-force extrapolation schedules are the same that were found in the optimization of Figures \ref{fig:overhead-plots} and \ref{fig:gate-depth}. \textit{Left:} We plot the operator norm error for every extrapolation schedule previously found. \textit{Right:} We introduce a layer of random single-qubit $X$-rotations to simulate measurement noise. The red dotted line signifies the noise rate corresponding to one standard deviation.}
    \label{fig:exact}
\end{figure}
